\centerline{\bf \Huge Учимся двигать}
\centerline{\bf \Huge Теория}

\begin{flushright}
        \scriptsize{}
\end{flushright}

\problem{0.1}
Найдите обратные преобразования к параллельному переносу, повороту, центральной и осевой симметрии и гомотетии.

\problem{0.2}
\textbf{Базовые повороты}
\begin{enumerate}
    \item Найдите поворот, переводящий квадрат в себя же
    \item Найдите поворот, переводящий правильный треугольник в себя же
    \item Найдите поворот, переводящий правильный $n$-угольник в себя же
    \item Две пересекающиеся прямые переводятся друг в друга поворотом.
\end{enumerate}

\problem{0.3}\textbf{Базовые осевые симметрии}
\begin{enumerate}
    \item Квадрат, окружность, равнобедренный треугольник и бесконечный угол имеют ось симметрии. Найдите их всех.
    \item$^*$ Композиция двух осевых симметрий является поворотом
\end{enumerate}


\problem{0.4}\textbf{Базовые гомотетии}
\begin{enumerate}
    \item Любые две окружности переводятся друг в друга гомотетией (вне зависимости от их взаимного расположения.)
    \item Касающиеся окружности переводятся друга в друга гомотетией с центром в точке касания
    \item Любые подобные треугольники переводятся друг в друга посредством поворотной гомотетии
    \item Любой треугольник гомотетичен своему серединному треугольнику
    \item Любые две параллельные прямые (даже если они совпадают) гомотетичны
    \item Любые два параллельных отрезка гомотетичны.
    \item Гомотетия переводит точку касания в точку касания
\end{enumerate}


\newpage

\hspace{3mm}

\centerline{\bf \Huge Учимся двигать}
\centerline{\bf \Huge Практика. Поворот}

\begin{flushright}
        \scriptsize{}
\end{flushright}

\problem{1}
Через центр квадрата проведены две перпендикулярные прямые. Докажите, что их точки пересечения со сторонами квадрата образуют квадрат.

\problem{2}
Докажите, что середины сторон правильного многоугольника образуют правильный многоугольник.

\problem{3}
    Точка $M$ лежит на дуге $AB$ описанной окружности правильного треугольника $ABC$. Докажите, что $MC = MA + MB.$\textit{(Теорема Помпею)}

\problem{4}
В ромбе $A B C D$ угол В равен $60^{\circ}$. На сторонах $BC$ и $CD$ выбраны точки $M$ и $N$ так, что $CM + CN = AB$. Докажите, что треугольник $AMN$ - правильный.

\problem{5}
$OABC$ и $ODEF$ --- квадраты (вершины идут по часовой стрелке). Найдите угол между прямыми $AD$ и $CF$.

\problem{6}
Два квадрата $BCDA$ и $BKMN$ имеют общую вершину $B$. Докажите, что медиана $B E$ треугольника $ABK$ и высота $BF$ треугольника $CBN$ лежат на одной прямой. (Вершины обоих квадратов перечислены по часовой стрелке.)

\problem{7}
На сторонах треугольника $A B C$ внешним образом построены правильные треугольники $A B C_1, A B_1 C$ и $A_1 B C$. Пусть $P$ и $Q$ - середины отрезков $A_1 B_1$ и $A_1 C_1$. Докажите, что треугольник $A P Q$ правильный.

\problem{8}
Точка Торриччели из курса геометрии.

\newpage

\hspace{3mm}

\centerline{\bf \Huge Учимся двигать}
\centerline{\bf \Huge Практика. Симметрия}

\begin{flushright}
        \scriptsize{}
\end{flushright}

\problem{1}
Четырёхугольник $ABCD$ таков, что $AB = CD$ , а $\angle \mathrm{BAD}=\angle \mathrm{CDA}$. Докажите, что $AD \parallel BC.$

\begin{wrapfigure}{r}{0.2\linewidth}
    \vspace{-7mm}
    \includegraphics[width=\linewidth]{Эльзята/GP_3.png}
\end{wrapfigure}

\problem{2}
Два квадрата и равнобедренный треугольник расположены так, как показано на рисунке. Докажите, что точки $\mathrm{A}, \mathrm{X}$ и $Y$ лежат на одной прямой.

\begin{wrapfigure}{r}{0.2\linewidth}
    \vspace{-7mm}
    \includegraphics[width=\linewidth]{Эльзята/GP_3_1.png}
\end{wrapfigure}

\problem{3}
Из точки $B$ к окружности проведены касательные $AB$ и $AC$ . Точки $P$ и $Q$ на окружности выбраны так, как показано на рисунке ( $\angle ABP =$ $\angle QBC$ ). Докажите, что отрезки $BP$ и $BQ$ равны.    

\problem{4}
В треугольнике $A B C$ угол $A$ равен $60^{\circ}$. Серединный перпендикуляр к отрезку $А B$ пересекает прямую $A C$ в точке $B_1$, а серединный перпендикуляр к отрезку $A C$ пересекает прямую $A B$ в точке $C_1$. Докажите, что прямая $B_1 C_1$ касается вписанной окружности треугольника $A B C$.

\problem{5}
Дан квадрат $ABCD$. На продолжении диагонали $AC$ за точку $C$ отмечена такая точка $K$ , что $BK = AC$. Найдите угол $BKC$.

\problem{6}
На стороне $CD$ ромба $A B C D$ нашлась такая точка $K$, что $AD = BK$. Пусть $F$ --- точка пересечения диагонали $BD$ и серединного перпендикуляра к стороне $ВС$. Докажите, что точки $A, F$ и $K$ лежат на одной прямой.

\problem{7}
Каждую вершину трапеции отразили симметрично относительно диагонали, не содержащей эту вершину. Докажите, что если получившиеся точки образуют четырехугольник, то он также является трапецией.

\problem{8}
Биссектрисы $AA_1$ и $CC_1$ треугольника $ABC$ пересекаются в точке $I$. Прямая, симметричная $АB$ относительно $CC_1$, пересекает прямую, симметричную $BC$ относительно $AA_1$, в точке $D$. Докажите, что $DI$ $\perp AC$.

\newpage

\hspace{3mm}

\centerline{\bf \Huge Учимся двигать}
\centerline{\bf \Huge Практика. Гомотетия.}

\problem{1}
С помощью задачи $0.4.(6)$ докажите \textit{замечательное свойство тpaпеции}: точка пересечения диагоналей, точка пересечения продолжений боковых сторон и середины оснований любой трапеции лежат на одной прямой.


\problem{2}
Даны две параллельные прямые $l_1$ и $l_2$. Окружности $\omega_1$ и $\omega_2$ касаются друг друга внешним образом в точке $B$, окружность $\omega_1$ касается $l_1$ в точке $A$ и $\omega_2$ касается $l_2$ в точке $C$. Докажите, что точки $A,B \text{ и } C$ лежат на одной прямой.

\problem{3}
Внутри квадрата $ABCD$ взята точка $M$. Докажите, что точки пересечения медиан треугольников $AMB, BMC, CMD, DMA$ образуют квадрат.

\problem{4}
На основаниях $BC$ и $AD$ трапеции $ABCD$ во внешнюю сторону построены равносторонние треугольники $BCK$ и $ADL$. Докажите, что $K,L$ и точка пересечения диагоналей лежат на одной прямой.

\problem{5}
Две окружности касаются в точке $К$. Прямая, проходящая через точку $К$, пересекает эти окружности в точках $A$ и $B$. Докажите, что касательные к окружностям, проведенные через точки $A$ и $B$, параллельны.

\problem{6}
Две окружности касаются в точке $K$. Через точку $K$ проведены две прямые, пересекающие первую окружность в точках $A$ и $B$, вторую - в точках $C$ и $D$. Докажите, что $A B \| C D$.

\problem{7}
На сторонах $A B$ и $A C$ треугольника $A B C$ взяты соответственно точки $M$ и $N$, причём $M N \| B C$. На отрезке МN взята точка $P$, причём $M P=\frac{1}{3} M N$. Прямая $A P$ пересекает сторону $B C$ в точке $Q$. Докажите, что $B Q=\frac{1}{3} B C$.

\problem{8}
а) Вписанная окружность треугольника $A B C$ касается стороны $A C$ в точке $D, D M$ - её диаметр. Прямая $B M$ пересекает сторону $A C$ в точке $K$. Докажите, что $A K=D C$.

\noindent
б) В окружности проведены перпендикулярные диаметры $A B$ и $C D$. Из точки $M$, лежащей вне окружности, проведены касательные к окружности, пересекающие прямую $A B$ в точках $E$ и $H$, а также прямые $M C$ и $M D$, пересекающие прямую $A B$ в точках $F$ и $K$. Докажите, что $E F=K H$.